\section{Related Work}
\label{sec:related}

Understanding why prior work did not measure the doctest feasibility gap requires examining
how code LLM data curation evolved --- from heuristic rules to quality-filtered corpora.

\subsection{Code LLM Benchmarks}

HumanEval~\cite{Chen2021Evaluating} established the pass@k metric for code generation evaluation,
using 164 hand-written Python problems with unit tests. MBPP~\cite{Austin2021Program}
provided 974 Python programming tasks ranging from simple to challenging. Both benchmarks
evaluate functional correctness via execution, making them the canonical measure of SFT
data quality effects.

\subsection{Heuristic-Only Corpus Filtering}

StarCoder~\cite{Li2023StarCoder} and StarCoder 2~\cite{Lozhkov2024StarCoder2} established
the standard pipeline for open-source code corpus construction from The Stack~\cite{Kocetkov2022Stack}:
near-deduplication, PII redaction, and heuristic quality filtering (line length, alphanumeric
fraction, encoding, XML/HTML ratio). These filters remove obvious junk but do not apply
execution-based quality signals. StarCoder achieves 40\% HumanEval pass@1 at 15.5B parameters.

\textbf{Limitation:} Heuristic filtering cannot distinguish syntactically valid from invalid
code or functionally correct from incorrect code.

\subsection{Quality Curation via Language Models}

phi-1~\cite{Gunasekar2023Textbooks} demonstrated that 7B quality tokens curated by GPT-4 yield
50.6\% HumanEval pass@1 at 1.3B parameters --- establishing that quality $>$ quantity at
equal token budget. phi-1.5~\cite{Li2023Textbooks2} extended this to reasoning tasks.

\textbf{Limitation:} GPT-4-based quality curation is not reproducible or scalable for
open-source use. Execution-based gates provide reproducible, cost-free quality signals.

\subsection{Execution-Based Data Selection}

EffiCoder~\cite{Zeng2024EffiCoder} applies execution-based sample selection to instruction-tuning
data (prompt-solution pairs), achieving +13pp HumanEval pass@1 on Qwen2.5-Coder-7B.
OpenCodeInstruct~\cite{OpenCodeInstruct2025} provides large-scale instruction-tuning data with execution
feedback. Token Cleaning~\cite{TokenCleaning2025} applies token-level influence filtering to SFT data.

\textbf{Limitation:} These works operate on instruction-tuning pairs with labeled test cases,
not raw corpus code. Raw corpus execution filtering requires intrinsic signals (\texttt{compile()},
doctest execution), and the feasibility of these signals at corpus scale has not been measured.

\subsection{Summary}

\begin{table}[h]
\centering
\caption{Comparison with prior work on code data curation.}
\label{tab:related}
\begin{tabular}{lllll}
\toprule
\textbf{Prior Work} & \textbf{Data Type} & \textbf{Execution Gate} & \textbf{Equal Budget} & \textbf{Feasibility Measured} \\
\midrule
StarCoder~\cite{Li2023StarCoder} & Raw corpus & Heuristic only & --- & No \\
phi-1~\cite{Gunasekar2023Textbooks} & GPT-4 curated & No & Yes & No \\
EffiCoder~\cite{Zeng2024EffiCoder} & Instruction pairs & Execution pass rate & No & No \\
\textbf{This work} & \textbf{Raw corpus} & \textbf{compile() + doctest} & \textbf{Yes (planned)} & \textbf{Yes} \\
\bottomrule
\end{tabular}
\end{table}
